\question[4]
Gegeben ist der folgende Python-Code:

\begin{lstlisting}
preis = 4.50                 # Zeile 1
anzahl = 3                   # Zeile 2
anzahl = anzahl + 2          # Zeile 3
gesamt = preis * anzahl      # Zeile 4
\end{lstlisting}

\begin{parts}
\part Gib alle Bezeichner an.
\part Gib für jede Zeile an, ob es sich um eine Initialisierung oder um eine (erneute) Wertzuweisung handelt.
\part Gib nach Zeile 4 für jede Variable ihren Wert und ihren Datentyp an.
\part Erläutere, warum Zeile 3 sinnvoll ist, obwohl sie als mathematische Gleichung keinen Sinn ergibt.
\end{parts}

\begin{solutionbox}{6cm}
\begin{parts}
\part \texttt{preis}, \texttt{anzahl}, \texttt{gesamt}
\part Zeile 1, 2 und 4: Initialisierung (erste Wertzuweisung). Zeile 3: erneute Wertzuweisung an \texttt{anzahl}.
\part \texttt{preis}: 4.5 (\texttt{float}), \texttt{anzahl}: 5 (\texttt{int}), \texttt{gesamt}: 22.5 (\texttt{float})
\part Eine Wertzuweisung wird von rechts nach links ausgeführt: Zuerst wird \texttt{anzahl + 2} mit dem alten Wert 3
ausgewertet, dann wird das Ergebnis 5 der Variable \texttt{anzahl} zugewiesen. Der alte Wert geht dabei verloren.
\end{parts}
\end{solutionbox}
